% Chapter introduction: why the topic arose, who first asked, what it answers.
% STATUS: written 2026-11-05 (sources read: Xiao 2015 pp. 1357--1368; Baclawski 1972
% pp. 351--356; Guo and Li 2014 pp. 545--572; Li and Wei 2018 pp. 79--92;
% Ortega 2006 pp. 225--243; Bennis and Fahid 2017 pp. 1--22; Boudine and Zerra 2024
% pp. 459--472; Lopatkin 2019 pp. 59--89).
Chapter~\ref{ch:matrix-triangular} read the derivations of algebras whose elements
are matrices, and its answers come from the shape of the matrix. An incidence
algebra keeps the shape of a matrix and replaces the entries by the intervals of a
poset \cite[p.~1358]{Xiao2015}. Ward first considered these algebras as generalized
algebras of arithmetic functions \cite[p.~1358]{Xiao2015}. Rota and Stanley then
made them the base of counting and of the theory of arithmetic functions
\cite[p.~1358]{Xiao2015}. This chapter asks for their derivations, and it
follows the same shape through a quiver, a trivial extension and a Leavitt path
algebra \cite[p.~59]{Lopatkin2019}.

Baclawski answered the derivation question for a locally finite ordered set
\cite[p.~352]{Baclawski1972}. Every derivation of the incidence algebra is an inner
derivation plus one attached to an additive function
\cite[Theorem~2, p.~352]{Baclawski1972}. When the poset has a smallest or a largest
element, the second part is inner as well \cite[p.~354]{Baclawski1972}. Xiao states
the same split for a pre-ordered set over a commutative ring, and he shows that no
Jordan derivation is new there \cite[Theorem~3.3, p.~1366]{Xiao2015}.

A path algebra gives a second reading of the same object \cite[p.~81]{LiWei2018}.
Guo and Li describe its derivations by conditions on the coefficients of an operator
on paths \cite[Theorem~2.3, p.~548]{GuoLi2014}. Their operators $D_{r,s}$ and the
inner derivations form a basis of the Lie algebra of derivations
\cite[Theorem~3.4, p.~555]{GuoLi2014}. Ortega shows that an incidence algebra is a
path algebra modulo the differences of parallel paths \cite[p.~232]{Ortega2006}. A
derivation of the incidence algebra is then a derivation of a path algebra that
kills the differences of parallel paths \cite[p.~232]{Ortega2006}.

A trivial extension adds a module with one crossed rule for the product
\cite[p.~4]{BennisFahid2017}. Bennis and Fahid describe its derivations in four
parts and decide which of them are inner
\cite[Theorem~3.12, p.~14]{BennisFahid2017}. A Leavitt path algebra adds a reverse
edge for every arrow \cite[p.~61]{Lopatkin2019}. Lopatkin describes its derivations,
and he shows that the outer derivations are infinite exactly when the graph has an
infinite path \cite[Corollary~3.8, p.~79]{Lopatkin2019}.

\hbfigure{11-path-incidence/poset-quiver}{A poset and its quiver. The incidence algebra of the poset is the path algebra of the quiver modulo the relation $\beta\alpha=\delta\gamma$.}{fig:path:poset}

Section~\ref{sec:path-incidence:path} gives the path algebra of a quiver. It states
the conditions of Guo and Li on a derivation, and the basis of the Lie algebra of
derivations \cite[Theorem~3.4, p.~555]{GuoLi2014}. It closes with the Jordan
derivations of a quiver with relations \cite[Theorem~3.4, p.~85]{LiWei2018}.
Section~\ref{sec:path-incidence:incidence} gives the incidence algebra, the two
theorems of Baclawski \cite[pp.~352, 354]{Baclawski1972} and the split of Xiao
\cite[Proposition~2.5, p.~1362]{Xiao2015}.
Section~\ref{sec:path-incidence:trivial-extension} gives the trivial extension and
the innerness criterion of Bennis and Fahid
\cite[Theorem~3.12, p.~14]{BennisFahid2017}.
Section~\ref{sec:path-incidence:leavitt} closes the chapter with the Leavitt path
algebra, its generators and the dimension of its outer derivations
\cite[Corollary~3.8, p.~79]{Lopatkin2019}.
